\section*{Chapter \ref{ch:g6:propdata} --- Proportionality and Data}

\begin{solution}{exo:g6:propdata:1}
Apples: \omterm{def:g3:division:remainder}{quotients} $\frac52 = 2.5$, $\frac{7.5}{3} = 2.5$,
$\frac{12.5}{5} = 2.5$ --- all equal: \omterm{def:g6:propdata:def}{proportional} (price
$2.50$ per kg).

Age/height: $\frac{85}{2} = 42.5$ but $\frac{103}{4} = 25.75$ --- not
equal: not \omterm{def:g6:propdata:def}{proportional}.
\end{solution}

\begin{solution}{exo:g6:propdata:2}
One croissant: $4.80 \div 4 = 1.20$ euros. Seven:
$7 \times 1.20 = 8.40$ euros.
\end{solution}

\begin{solution}{exo:g6:propdata:3}
Coefficient: $18 \div 10 = 1.8$ euros per liter. Then
$25$ L cost $25 \times 1.8 = 45$ euros; $72$ euros buy
$72 \div 1.8 = 40$ L; $60$ L cost $60 \times 1.8 = 108$ euros.
\end{solution}

\begin{solution}{exo:g6:propdata:4}
For $250$ km: $2.5$ times the fuel of $100$ km, so
$2.5 \times 6 = 15$ L. For $350$ km: $3.5 \times 6 = 21$ L. With
$27$ L: $27 \div 6 = 4.5$ hundreds of km, i.e.\ $450$ km.
\end{solution}

\begin{solution}{exo:g6:propdata:5}
$50\,\%$ of $84$: \omterm{def:g3:division:half}{half}, $42$. \quad
$25\,\%$ of $200$: a \omterm{def:g3:division:half}{quarter}, $50$. \quad
$10\,\%$ of $63$: a tenth, $6.3$. \quad
$30\,\%$ of $70$: $3 \times 7 = 21$.
\end{solution}

\begin{solution}{exo:g6:propdata:6}
$60\,\%$ of $450$: $\frac{60}{100} \times 450 = 270$ students eat at
the cafeteria; $450 - 270 = 180$ do not.
\end{solution}

\begin{solution}{exo:g6:propdata:7}
Fewer than $10$ books: Tuesday ($8$) and Thursday ($6$). Wednesday
minus Thursday: $17 - 6 = 11$ more books.
\end{solution}

\begin{solution}{exo:g6:propdata:8}
Bar chart with five bars of heights $14$, $16$, $13$, $17$, $20$
(graduation every $2$ or $5$ degrees works well). Warmest day: Friday
($20$ degrees).
\end{solution}

\begin{solution}{exo:g6:propdata:9}
For $10$ people, multiply by $\frac{10}{6} = \frac53$: flour
$450 \times \frac{10}{6} = 750$ g. Eggs:
$3 \times \frac{10}{6} = 5$ --- luckily a whole number. (If it were
not, one would \omterm{def:g4:numbers:round}{round} \emph{up} to have enough.)
\end{solution}

\begin{solution}{exo:g6:propdata:10}
\emph{1.} In $2$ h: $48$ km. In \omterm{def:g3:division:half}{half} an hour: $12$ km. In
$1$ h $30$: $24 + 12 = 36$ km.

\emph{2.} $60 \div 24 = 2.5$ hours, i.e.\ $2$ h $30$ min.
\end{solution}

\begin{solution}{exo:g6:propdata:11}
Ball: discount $20\,\%$ of $15 = 3$ euros, new price $12$ euros.
Racket: discount $8$ euros, new price $32$ euros. The new price is
$80\,\%$ of the old one in every case: \omterm{def:g6:propdata:def}{proportional}, coefficient
$0.8$.
\end{solution}

\begin{solution}{exo:g6:propdata:12}
After $t$ hours, the thick candle has burned $\frac t6$ of its height:
it stands at $1 - \frac t6$ of the initial height; the thin one at
$1 - \frac t4$. We want
\[
1 - \frac t4 = \frac12 \left(1 - \frac t6\right)
\quad\text{i.e.}\quad
1 - \frac t4 = \frac12 - \frac t{12}.
\]
Then $\frac12 = \frac t4 - \frac t{12} = \frac{3t - t}{12} =
\frac{t}{6}$, so $t = 3$ hours. Check: after $3$ h the thick candle
stands at $1 - \frac36 = \frac12$ and the thin one at
$1 - \frac34 = \frac14$ --- indeed \omterm{def:g3:division:half}{half} of it.
\end{solution}

\begin{solution}{pb:g6:propdata:1}
\textbf{1.} $100\,000$ cm $= 1\,000$ m $= 1$ km: one centimeter
on the map is one kilometer on the ground.

\textbf{2.} $7.5$ cm $\rightarrow 7.5$ km between the villages;
$12$ cm $\rightarrow 12$ km of lake shore.

\textbf{3.} $20$ km $\rightarrow 20$ cm on the map.

\textbf{4.} $5$ km $= 500\,000$ cm, so the scale is
$1:500\,000$. The hiking map ($1:100\,000$) is the more
detailed: it uses $5$ cm of paper where the other spends only
$1$ cm.

\textbf{5.}
\begin{center}
\renewcommand{\arraystretch}{1.2}
\begin{tabular}{l|ccc}
map (cm) & $1$ & $2$ & $7.5$ \\
\hline
real (km) & $1$ & $2$ & $7.5$ \\
\end{tabular}
\end{center}
Real distance $=$ map distance $\times 1$ (in these units):
always the \emph{same} multiplier, which is exactly the
definition of \omterm{def:g6:propdata:def}{proportionality}
(\cref{def:g6:propdata:def}). The coefficient here is $1$
kilometer per centimeter.

\textbf{6.} $4$ m $= 400$ cm and $400 \div 50 = 8$;
$3$ m $= 300$ cm and $300 \div 50 = 6$: the plan shows an
$8$ cm $\times$ $6$ cm rectangle.

\textbf{7.} Bed: $200 \div 50 = 4$ cm by $90 \div 50 = 1.8$ cm.
Doorway: $80 \div 50 = 1.6$ cm.

\textbf{8.} Real \omterm{def:g6:measure:area}{area}: $4 \times 3 = 12$ m$^2$. Plan \omterm{def:g6:measure:area}{area}:
$8 \times 6 = 48$ cm$^2$. Converting:
$12$ m$^2 = 12 \times 10\,000 = 120\,000$ cm$^2$, and
\[
120\,000 \div 48 = 2\,500 .
\]
The \omterm{def:g3:division:remainder}{quotient} is not $50$ but $2\,500 = 50 \times 50$.

\textbf{9.} One cm$^2$ of paper stands for a real square of
$50$ cm by $50$ cm, which contains
$50 \times 50 = 2\,500$ cm$^2$. So when lengths are divided by
$50$, \omterm{def:g6:measure:area}{areas} are divided by $50 \times 50 = 2\,500$: \omterm{def:g6:measure:area}{areas} follow
the \emph{square} of the scale.

\textbf{10.} $6 \times 2\,500 = 15\,000$ cm$^2$, and
$15\,000 \div 10\,000 = 1.5$ m$^2$.

\textbf{11.} The picture suggests Sunday sold \emph{twice} as
much (a bar twice as tall). Truth: the increase is $5$ euros out
of $105$, and $5 \div 105 \approx 0.048$: about $5\,\%$ more,
not $100\,\%$ more. The manager ignored the warning to check
that the axis starts at $0$
(\cref{met:g6:propdata:read}).

\textbf{12.} With the axis from $0$, the bars rise to $105$ and
$110$ small units: two bars of nearly the same height, the
honest picture of a $5\,\%$ \omterm{ex:g1:subtraction:difference}{difference}.

\textbf{13.} Up: the increase is $20$ stamps from a start of
$40$: $\frac{20}{40} = 50\,\%$. Down: the decrease is $20$
stamps from a start of $60$: $\frac{20}{60} = \frac13 \approx
33\,\%$. Same $20$ stamps, different starting values --- a
percentage is always a \omterm{def:g6:fractions:def}{fraction} \emph{of something}, and the
``something'' changed.

\textbf{14.} After $20\,\%$ off: $100 - 20 = 80$ euros. The
extra $10\,\%$ applies to $80$: discount $8$ euros, final price
$72$ euros. Total discount: $28$ euros out of $100$, so
$28\,\%$ --- not $30\,\%$. The second discount acted on the
already-reduced price, so its euros are smaller: percentages of
different quantities do not add.

\textbf{15.} On that map, $1$ cm stands for $5$ km, so $1$
cm$^2$ stands for $5 \times 5 = 25$ km$^2$ (question 9's rule).
The forest: $3 \times 25 = 75$ km$^2$. Lengths follow the
scale; \omterm{def:g6:measure:area}{areas} follow its \emph{square}.
\end{solution}
